\section{Empirical Strategy}\label{sec:strategy}

\subsection{Stacked difference-in-differences}

The treatment---absorbing ex-cartel workers---is staggered: different firms absorb workers in different years (2012--2015), as cartel conduct periods end at different times. A standard two-way fixed-effects (TWFE) specification is subject to the biases documented by \citet{goodman2021difference}: already-treated units serve as implicit controls for later-treated units, potentially contaminating the estimates. We adopt the stacked-cohort approach of \citet{cengiz2019effect} to avoid this.

For each cohort $g \in \{2012, 2014, 2015\}$, we construct a separate clean dataset:\footnote{We drop cohort 2013 ($N = 34$ treated firms) because its small size generates noisy pre-trend estimates that do not permit meaningful inference. Results including cohort 2013 are qualitatively similar but less precise.}

\begin{itemize}
\item \textbf{Treated}: firms whose first ex-cartel worker arrived in year $g$.
\item \textbf{Controls}: never-treated firms (firms that bid on cartel item codes but never absorbed an ex-cartel worker).
\item \textbf{Window}: $[g - 2, g + 4]$.
\end{itemize}

The cohort-specific datasets are stacked, and we estimate:
\begin{equation}\label{eq:stacked}
y_{ijt}^{(g)} \;=\; \alpha_i^{(g)} + \gamma_{jt}^{(g)} + \sum_{\tau \neq -1} \beta_\tau \cdot D_{i\tau}^{(g)} + \varepsilon_{ijt}^{(g)}\,,
\end{equation}
where $i$ indexes firms, $j$ indexes item codes, $t$ indexes calendar years, $g$ indexes cohorts, $\alpha_i^{(g)}$ is a firm-within-cohort fixed effect, $\gamma_{jt}^{(g)}$ is an item-code-by-year-within-cohort fixed effect, and $D_{i\tau}^{(g)} = 1$ if firm $i$ is treated in cohort $g$ and the observation is at event time $\tau = t - g$. The omitted category is $\tau = -1$. Standard errors are clustered at the firm-within-cohort level.

The item-code-by-year-within-cohort fixed effect $\gamma_{jt}^{(g)}$ absorbs all market-level variation---changes in demand, reference prices, and the number of competitors---within each cohort. The identifying assumption is parallel trends: absent treatment, the treated and never-treated firms would have followed the same within-market trajectory. We assess this by inspecting pre-treatment coefficients.

\subsection{Endogeneity of treatment timing}\label{sec:endogeneity}

Equation~\eqref{eq:stacked} treats the year a firm first absorbs an
ex-cartel worker as a shock to its competitive behavior. Hiring, however,
is a choice, and the timing of cartel-worker arrival is plausibly
correlated with unobserved determinants of the outcome. We confront this
concern through three complementary strategies: (i)~a shift-share exposure
design that instruments worker absorption with predetermined occupational
and geographic overlap between the receiving firm and the cartel; (ii)~a
conditional-on-hiring placebo that holds the firm's hiring decision fixed
and varies only the source of the new worker; and (iii)~the imputation
estimator of \citet{borusyak2024revisiting} as a check against
staggered-treatment artifacts.

Four channels of endogeneity are conceptually distinct, and all bias the
headline coefficient in the same direction. \emph{First}, capacity-pull
reverse causality: a firm $i$ that wins more cartel-item contracts at
date~$t$ requires labor to deliver them, posts vacancies, and absorbs
workers from a cartel firm that is simultaneously contracting---so the
arrival of the worker may mark $i$'s expansion rather than cause it.
\emph{Second}, common demand shocks: a buyer-side expansion in product
class~$j$ can raise both $i$'s win rate in~$j$ and $i$'s demand for
$j$-relevant labor, generating correlation through a third variable that
the item-by-year fixed effect~$\gamma_{jt}$ does not absorb when the shock
loads on a subset of firms. \emph{Third}, mechanical displacement on
controls: in zero-sum auctions, never-treated firms competing on the same
items lose contracts when the treated firm wins them, so the
difference-in-differences estimator captures the treatment effect plus a
displacement term. \emph{Fourth}, network selection: workers move along
occupational and geographic ties, and firms with overlapping labor pools
tend to have overlapping product portfolios, so the treatment indicator
selects firms that are close substitutes for the cartel firm in the
product market---precisely the firms whose win rates rise mechanically as
the cartel exits. All four channels push~$\hat{\beta}$ upward; we
therefore interpret the OLS stacked-DiD estimates as upper bounds on the
true capacity-transfer effect.

\subsubsection*{Shift-share exposure instrument}

To address reverse causality and common demand shocks, we construct a
shift-share instrument in the spirit of \citet{borusyak2025nonrandom}. For
each receiving firm~$i$ and year~$t$, define
\begin{equation}\label{eq:bartik}
Z_{it} \;=\; \sum_{c \in \mathcal{C}} \mathbb{1}[t \geq T_c]
            \cdot W_c \cdot s_{i,c}^{\text{pre}}\,,
\end{equation}
where $\mathcal{C}$ is the set of convicted cartels, $T_c$ is the
conduct-end date of cartel~$c$, $W_c$ is the cartel's end-of-conduct
workforce by occupation--municipality cell, and $s_{i,c}^{\text{pre}}$ is
firm~$i$'s pre-period (2009) employment share in the same
occupation--municipality cells. The instrument captures the structural
exposure of firm~$i$ to the labor reallocation triggered by enforcement,
using only the cartel's contraction (the ``shift,'' exogenous to~$i$
conditional on enforcement timing) and $i$'s predetermined
occupational--geographic mix (the ``share,'' fixed before any cartel
worker could have moved).

We use~$Z_{it}$ as an instrument for the number of ex-cartel workers
absorbed by firm~$i$ in year~$t$ in a 2SLS event-study analogue of
equation~\eqref{eq:stacked}. The exclusion restriction is that
firm~$i$'s 2009 occupational--geographic structure is uncorrelated with
its post-conviction idiosyncratic demand shocks in cartel item codes,
conditional on the item-by-year fixed effect. We assess this restriction
by regressing $i$'s \emph{pre}-period item-level win-rate trajectories on
lagged~$Z_{it}$; the absence of predictive power supports the
exclusion.\footnote{Following \citet{borusyak2025nonrandom}, identification
requires either many quasi-random shocks or exogenous shares. With only
six cartels contributing shocks, we are at the lower end of what the
asymptotic theory supports, and we treat the IV as a complement to---not a
substitute for---the OLS upper bounds.}

\subsubsection*{Conditional-on-hiring placebo}

The shift-share design addresses confounders that operate through the
hiring decision; a complementary test holds the hiring decision fixed and
varies only the source of the new worker. We restrict the sample to
firm-years in which firm~$i$ hired at least one worker in an occupation
that overlaps with the cartel-firm workforce, and estimate
\begin{equation}\label{eq:condhire}
y_{ijt} \;=\; \alpha_i + \gamma_{jt}
            + \beta \cdot \text{CartelHire}_{it}
            + \delta \cdot \text{AnyHire}_{it}
            + \varepsilon_{ijt}\,,
\end{equation}
so that~$\beta$ is identified from variation in the \emph{source} of the
new worker rather than from the firm's decision to hire. A null on~$\beta$
would indicate that the headline effect reflects the hiring decision
itself; a surviving positive~$\beta$ supports the interpretation of the
worker as a vehicle of cartel-firm-specific knowledge or capacity.

\subsubsection*{Robust estimator and imputation}

We additionally report the imputation estimator of
\citet{borusyak2024revisiting}, which uses both never-treated and
not-yet-treated firms to construct counterfactuals and which has correct
standard errors under heterogeneous treatment effects. Divergence between
the imputation estimator and the stacked DiD would localize residual
staggered-treatment bias. Results for all three robustness strategies are
reported in Section~\ref{sec:rob_endog}.

\subsection{CEM matching}

As a robustness check, we implement Coarsened Exact Matching (CEM) on pre-period characteristics measured in the two years before the first ex-cartel worker arrives: the number of bids, the number of distinct item codes, the win rate, and firm size (number of RAIS employees). We match on quintiles of each variable and estimate a weighted TWFE event study on the matched sample, with CEM weights applied. This addresses the concern that treated and control firms differ systematically in pre-period levels, which could generate spurious pre-trends.

\subsection{Exit placebo}

To test whether the treatment effect is cartel-specific, we construct a placebo by replacing the cartel-worker treatment with an exit-worker treatment: absorption of workers from non-cartel BEC firms that stopped bidding between 2011 and 2014. If the win-rate and price effects are driven by cartel-specific knowledge transfer, they should be absent---or at least substantially attenuated---in the exit placebo. If they are similar, the effect is generic to any labor-supply shock from a contracting firm.
